% sections/introduccion.tex
\section*{Introducción}
La ciberseguridad busca proteger la información y los sistemas frente a diferentes amenazas que pueden comprometer su confidencialidad, integridad y disponibilidad. En este taller se analizan casos reales que permiten comprender cada uno de estos principios: la fuga de información de Ashley Madison como ejemplo de afectación a la confidencialidad, la modificación no autorizada del software CCleaner como caso de pérdida de integridad y el ataque a Shinhan Bank y otras instituciones de Corea del Sur como ejemplo de afectación a la disponibilidad.

Finalmente, se abordan los conceptos de Autenticación, Autorización y Accounting (AAA) y se relacionan con los principales patrones de brechas de seguridad identificados en el 2025 Verizon Data Breach Investigations Report, con el propósito de entender cómo estos mecanismos contribuyen a prevenir, detectar y limitar los efectos de los incidentes de ciberseguridad.
